% ID: FIS-ROT-ANG-002
% Titolo: Otto studenti su una giostra rotante
% Disciplina: Fisica
% Area: Dinamica rotazionale
% Argomento: Conservazione del momento angolare
% Sottoargomento: Variazione del momento d'inerzia in un sistema ridotto
% Classe: Quarta liceo scientifico
% Difficolta: 3
% Tipo: problema_numerico
% Risultato: due studenti; omega_f circa 1,21 rad/s
% Tempo_stimato: 10 min
% Competenze: modellizzazione, conservazione del momento angolare, momento d'inerzia
% Prerequisiti: moto circolare, momento d'inerzia, velocita angolare
% Tag: fisica, dinamica_rotazionale, momento_angolare, momento_inerzia, sistema_ridotto, simulazione
% Fonte: esercizio_originale_exergo
% Autore: Riccardo Carli
% Licenza: CC-BY-SA-4.0
% Simulazione: rotational_platform

\begin{esercizio}
Una giostra circolare, modellizzata come un disco pieno, ha raggio
\(1{,}50\,\mathrm{m}\) e massa \(240\,\mathrm{kg}\). Ruota senza attrito
attorno al proprio asse con velocita angolare iniziale
\[
\omega_i=1{,}00\,\mathrm{rad/s}.
\]
Otto studenti di massa \(50{,}0\,\mathrm{kg}\) ciascuno si trovano,
equidistanti, alla distanza \(1{,}20\,\mathrm{m}\) dall'asse.

Nel modello didattico dell'esercizio si considera costante il momento angolare
di riferimento del sistema ridotto formato dalla giostra e dagli studenti che
vi restano. Ogni studente e modellizzato come una massa puntiforme. Quanti
studenti devono scendere dalla giostra affinche la velocita angolare diventi
circa \(1{,}20\,\mathrm{rad/s}\)?

Nota sul modello: uno studente che lascia realmente la giostra puo portare con
se momento angolare. L'ipotesi di conservazione applicata qui al sistema
ridotto e una semplificazione didattica, non una descrizione completa della
fase di uscita.
\end{esercizio}

\begin{soluzione}
Il momento d'inerzia della giostra vale
\[
I_\text{giostra}=\frac{1}{2}MR^2
=\frac{1}{2}\cdot240\cdot(1{,}50)^2
=270\,\mathrm{kg\,m^2}.
\]

Ogni studente contribuisce con
\[
I_\text{studente}=mr^2
=50{,}0\cdot(1{,}20)^2
=72\,\mathrm{kg\,m^2}.
\]

Con otto studenti il momento d'inerzia iniziale e
\[
I_i=270+8\cdot72
=846\,\mathrm{kg\,m^2},
\]
e il momento angolare di riferimento assunto dal modello e
\[
L_\text{rif}=I_i\omega_i
=846\cdot1{,}00
=846\,\mathrm{kg\,m^2/s}.
\]

Se scendono \(k\) studenti, ne restano \(8-k\) e
\[
\omega_f(k)=\frac{L_\text{rif}}
{270+(8-k)\cdot72}.
\]

Con due studenti in meno ne restano sei:
\[
I_f=270+6\cdot72
=702\,\mathrm{kg\,m^2},
\]
\[
\omega_f=\frac{846}{702}
\simeq1{,}205\,\mathrm{rad/s}.
\]
Il valore e circa \(1{,}20\,\mathrm{rad/s}\). Devono quindi scendere dalla
giostra \textbf{due studenti}, nel modello didattico dichiarato dal testo.
\end{soluzione}
